\chapter{Deep Research 領域深度研究}

\section{領域知識與核心設計觀點}

餐廳推薦涉及使用者對用餐情境的主觀偏好，公車旅程則受到班次、步行與轉乘時序約束。兩者整合時，須同時處理偏好表達、可行性判定及資料證據的差異。以下由相關文獻歸納五項設計觀點，作為第四章需求表示與推薦流程的依據。

\subsection{必要限制與偏好效用的分離}

形式化規劃以限制條件界定可行解，主觀屬性研究則指出使用者對描述語句的理解與偏好可能不同~\cite{llm_planning_tools,subjective_attributes}。兩類資訊在決策中的作用因而有所區別：預算與轉乘時序限制決定方案是否可接受，聚餐氛圍等偏好則用於比較可接受方案。若僅將兩者合併為加權分數，較高的語意相關性仍可能抵銷已知違反必要條件的成本。

本研究據此將可行性查核與偏好排序分離，由程式判定必要限制，再對通過者計算情境與交通分數。使用者若將主觀偏好指定為必要條件，系統亦須記錄其判定標準與支持證據，避免僅憑相似度分數宣稱符合。

\subsection{交通與營業條件的時間相依性}

公共運輸規劃以有先後順序的班次連接與步行關係描述旅程；餐廳推薦研究亦已將請求時間與地理位置納入評分~\cite{bus_dynamic_replanning,restaurant_habib}。對結合兩者的推薦問題而言，餐廳是否適合前往，取決於完成候車、乘車與轉乘後的實際抵達時刻。因此，查詢當下的營業狀態不足以單獨決定用餐可行性。

本研究以出發、上車、轉乘、抵達及離店事件構成完整時間線，並使用抵達至離店區間查核營業條件。候車或轉乘時間更新時，後續事件須依相同時間基準重新計算；營業與步行條件的完整判定亦須在取得交通方案後完成。

\subsection{預測不確定性與證據完備性}

Rodriguez-Deniz 與 Villani 以機率分布建模公車延誤，顯示到站預測的評估需要考慮點估計以外的不確定性~\cite{bus_delay_uncertainty}。另一方面，資料缺漏或班次無法對齊，屬於系統欠缺足夠判定依據的情形。預測誤差與證據不足雖均可能影響推薦可靠性，兩者的表示與處理方式並不相同。

本研究以三值查核表示證據下的通過、不通過與不確定，並保留外部預估值、參考時刻及緩衝設定。三值狀態描述目前能否支持條件判定，不具有機率意義；查核通過亦僅表示方案在所採資料與規則下成立。

\subsection{推薦推論的證據可追溯性}

評論式對話推薦透過文字片段支持偏好判斷，知識圖譜增強推薦則利用結構化資訊協助模型產生結果~\cite{critique_preference,knowledge_graph_llm_recommendation}。整合此類資訊時，店家事實、評論中的主觀描述與系統推算應保留各自的來源，才能追查推薦理由所依據的資料及計算過程。

因此，本研究分別保存來源的發布、觀測、取得與適用時間，並將推薦說明連結至條件判定紀錄。影片發布日期僅描述來源的發布事件，不能據此認定其中價格或營業資訊持續有效；圖譜關係則用於組織證據，不直接作為資訊正確性的保證。

\subsection{目的地選擇與交通方案的相依關係}

時空地點推薦研究將位置與到訪時間納入偏好建模，公共運輸研究則分析班次與轉乘對旅程的影響~\cite{poi_stisan,bus_dynamic_replanning}。由此可見，同一間餐廳在不同出發地與時刻下，可能具有不同的交通負擔；目的地的適合程度不能完全獨立於到達方式評估。

本研究先以評論語意形成餐廳候選集合，再分析各餐廳可用的公車方案，將交通查核與便利性納入最終排序。同一餐廳僅保留其合適方案參與餐廳間比較，以避免路線數量較多者重複占據推薦結果。此流程將搜尋範圍限定於情境相符的候選，並以交通條件進一步決定取捨。

\section{關鍵方法取捨與研究議題}

\subsection{時刻表規劃與即時資訊的整合}

時刻表描述服務日、班次及預定連接關係，提供完整旅程規劃的基礎。即時資訊則反映可觀測的延誤與服務變動，能修正出發前規劃與當下狀態之間的差異。Abuaisha 等人的動態重規劃方法以時刻表網路為基礎，比較靜態規劃、出發前快照及行程中更新，並以 pull 與 push 機制分析更新所需的查詢與計算負擔~\cite{bus_dynamic_replanning}。因此，兩類資料具有互補作用，其整合重點在於更新範圍與成本的取捨。

該研究使用四個城市的 GTFS 網路，但延誤來自合成模型，真實延誤資料的驗證仍列為後續研究~\cite{bus_dynamic_replanning}。這項結果支持即時更新具有研究價值，卻不足以直接決定臺中場域的資料覆蓋、更新間隔或查詢頻率。本研究因而以班表與旅運規劃形成候選，於可取得即時依據的範圍內查核搭乘條件，並保存所採時間基準。行程中持續監控另列為擴充，以維持初版查詢流程的明確範圍。

\subsection{旅程時間與轉乘可靠性的取捨}

最早抵達提供明確的路線比較目標，動態重規劃亦以抵達時間衡量更新方案的效益~\cite{bus_dynamic_replanning}。然而，若方案的轉乘餘裕過小，僅依單一預估時間排序，可能無法充分反映延誤對後續搭乘的影響。Rodriguez-Deniz 與 Villani 對公車延誤分布的分析，則從預測模型層次說明考慮不確定性的必要性~\cite{bus_delay_uncertainty}。

兩類研究分別關注路線決策與預測分布，並非同一實驗下對兩種策略的直接比較；其證據尚不足以建立適用所有路線與使用者的固定緩衝值。增加緩衝可以保留較大的銜接餘裕，亦可能排除原本可搭乘的方案，或增加旅程時間。因此，本研究將緩衝視為需依場域資料與需求設定的參數，並保存其版本；時間模型呈現候車與轉乘條件，推薦說明則揭露所採預估與緩衝，不將查核狀態解讀為實際成功機率。

\subsection{語言模型規劃與程式化查核的分工}

語言模型能以自然語言分解任務與修訂計畫。HyperTree Planning 透過階層式規劃結構改善模型推理，顯示模型本身的規劃能力仍具有可發展空間~\cite{hypertree_planning}。另一條研究路徑則將需求轉為形式化條件，由求解工具產生方案及診斷不可滿足的限制。Hao 等人的研究在其設定下取得較高的限制滿足率，同時指出問題編碼、求解時間與外部資料正確性仍需處理~\cite{llm_planning_tools}。

上述方法的任務、模型與評估設定不同，不能藉由跨論文的成功率直接判定何者在本場域較佳。本研究依條件可追溯與規則一致的需求，採用語言模型整理輸入、程式建立時間線與查核限制的分工。此選擇仍仰賴需求解析的正確性：即使後續程式正確，入口的預算單位或日期若解析錯誤，系統仍可能驗證不符合使用者原意的條件。因此，解析欄位須保留與原始語句的對應，並對影響必要限制的歧義進行補問。

\section{領域知識的推理檢核}

表~\ref{table:deep_questions} 以十項情境問題檢核上述觀點，涵蓋時間相依性、證據衝突與方法取捨。各題著重能否說明判定依據、推導失敗情境及辨識推論範圍，而非僅辨認技術名詞。

{\small\singlespacing
\renewcommand{\arraystretch}{1.25}
\setlength{\tabcolsep}{3pt}
\begin{longtable}{|P{\dimexpr 0.1500\linewidth-2\tabcolsep-1.3333\arrayrulewidth\relax}|P{\dimexpr 0.4600\linewidth-2\tabcolsep-1.3333\arrayrulewidth\relax}|P{\dimexpr 0.3900\linewidth-2\tabcolsep-1.3333\arrayrulewidth\relax}|}
\caption{十個深度理解問題與推理重點}\label{table:deep_questions}\\
\hline
\textbf{類型} & \textbf{問題} & \textbf{應能推理的重點} \\ \hline
\endfirsthead
\hline
\textbf{類型} & \textbf{問題} & \textbf{應能推理的重點} \\ \hline
\endhead
Why & 為何營業查核要使用抵達與離店時間，而不是查詢時間？ & 候車與轉乘改變後續時間；需覆蓋用餐區間 \\ \hline
Failure case & 預估 18:50 抵達、19:00 打烊，用餐需 90 分鐘，應如何判定？ & 抵達時有開門仍不足；用餐結束超過打烊為不通過；數字為設計示例 \\ \hline
What-if & 公車預估 3 分鐘到站，使用者走到站牌需 6 分鐘，且這是最後一班時，能否推薦？ & 正常狀態與 ETA 存在不代表能搭上；須比較準備時間與班次 \\ \hline
Trade-off & 增加轉乘緩衝後錯誤減少，但已驗證覆蓋也大幅下降，能否稱為改善？ & 同時報告風險與覆蓋，並比較旅程成本 \\ \hline
Evidence & 兩個來源對營業時間不同，應依什麼資訊決定可用證據？ & 來源、適用日期、核對時間、版本與衝突；不能單看抓取時間 \\ \hline
Comparison & 若 TDX 與另一服務的 ETA 不同，如何判斷誰較準？ & 同站位、方向、班次、觀測時點對齊，另取實際到站紀錄 \\ \hline
Design decision & 圖譜與 SQL 保存相同事實時，圖譜仍可能有什麼價值？ & 證據路徑、查詢正確性、延遲與維護成本；不能預設推薦一定更好 \\ \hline
Real-world scenario & 已上車後轉乘取消，為何不能仍以原出發地規劃？ & 行程狀態、當前可下車站與新時間；此為後續動態擴充應處理的情境 \\ \hline
Counterexample & 所有查詢都不回傳方案，限制違反率很低，是否代表可靠？ & 零分母不能記為零錯誤；須看拒答、覆蓋與假性無解 \\ \hline
Evaluation & 只重播 N1 紀錄，能否證明推薦降低實際錯過公車的比例？ & 可以驗證資訊變動與處理流程；實際到站與反事實效果須有額外證據 \\ \hline
\end{longtable}
}

